\documentclass[14pt]{extarticle}
\usepackage{fontspec}
\usepackage{polyglossia}
\usepackage{amsmath}
\usepackage{amsfonts}
\usepackage{amssymb}
\usepackage{pgfplots}
\pgfplotsset{compat=1.18}

% Hebrew setup
\setdefaultlanguage{hebrew}
\setotherlanguage{english}
\newfontfamily\hebrewfont[Script=Hebrew]{Arial}

% Math font setup for proper Greek symbols
\usepackage{unicode-math}
\setmathfont{Latin Modern Math}

% Increase line spacing throughout document
\linespread{1.3}

\usepackage[margin=1in]{geometry}

% Left-aligned equation environment
\newenvironment{lefteqs}{%
    \begin{english}%
    \begin{flushleft}%
    \renewcommand{\arraystretch}{1.3}%
    $\begin{array}{@{}l@{}}%
}{%
    \end{array}$%
    \end{flushleft}%
    \end{english}%
}

\begin{document}

\begin{center}
    {\LARGE \textbf{מטלה 4 - אלקטרומגנטיות אנליטית}}\\[0.5em]
    {\large שי רואימי - 318986270}\\[0.3em]
\end{center}

\hrule
\vspace{2em}

\textbf{\large שאלה 1}

(פרק 4, שאלה 6)

נתון שדה מגנטי תלוי זמן:
\begin{lefteqs}
    \vec{H}(t) = H_0(t)\hat{x} + H_1(t)\hat{y}
\end{lefteqs}

נמצא את הסירקוליציה של השדה החשמלי על המסילה שבציור.

נשתמש בחוק פאראדיי האינטגרלי:
\begin{lefteqs}
    \oint_C \vec{E}\cdot d\vec{\ell} = -\frac{d}{dt}\iint_S \vec{B}\cdot \hat{n}\,dS, \qquad \vec{B}=\mu_0\vec{H}
\end{lefteqs}

כיוון הנורמל לשטח הוא $\hat{n}=\hat{y}$, ע"פ חוק יד ימין
לבחירת נורמל $\hat{n}=\hat{y}$ עבור הלולאה הנתונה מתקבל:
\begin{lefteqs}
    -\frac{d}{dt}\iint_S \mu_0\left(H_0(t)\hat{x}+H_1(t)\hat{y}\right)\cdot\hat{y}\,dS
    = -\frac{d}{dt}\left(\mu_0\iint_S H_1(t)\,dS\right)
    = \mu_0wl\,\frac{dH_1}{dt}
\end{lefteqs}

ולכן:
\begin{lefteqs}
\oint_C \vec{E}\cdot d\vec{\ell} = \mu_0wl\,\frac{dH_1}{dt}
\end{lefteqs}

\vspace{3em}
\hrule
\vspace{2em}

\textbf{\large שאלה 2}

(פרק 4, שאלה 8, סעיף ה')

נתון מעגל זרם מתחלף בעל מתח מקור:
\begin{lefteqs}
V(t)=V_0e^{-\gamma t}
\end{lefteqs}

במעגל ישנו קבל (שני לוחות בעלי שטח S, רדיוס a) במרחק $d$ אחד מהשני.

גליל באורך $d$ ורדיוס $b$ ($b<a$) עשוי חומר מוליף בעל מוליכות $\sigma$ ומקדם דיאלקטרי $\epsilon$ בעוד ההיקף החיצוני ריק.

\textbf{א)} 
לכל $0<r<a$, מתקיים $\vec{E}=\frac{V(t)}{d}\hat{z}=\frac{V_0e^{-\gamma t}}{d}\hat{z}$

מתח המקור מקבע את הפרש הפוטנציאלים בין הלוחות.

\textbf{ב)} 
נעזר בחוק אמפר המתוקן

\begin{lefteqs}
    \oint_C\vec{B}\cdot d\vec{\ell}=\mu_0I_{in}+\mu_0\epsilon_0\frac{\partial\Phi_E}{\partial t}
\end{lefteqs}

כאשר עבור $r<b$, התווך הוא הגליל המוליך, נסתכל על לולאה אמפירית באזור בין שני הלוחות בריוס $r$:



\begin{lefteqs}
    I_{in}=J \cdot S = \sigma\pi r^2E = \sigma\pi r^2 \frac{V_0e^{-\gamma t}}{d} \\
    \Phi_E = E \cdot S = \frac{V_0e^{-\gamma t}}{d} \cdot \pi r^2 \\
    \frac{\partial\Phi_E}{\partial t}=-\gamma\pi r^2 \frac{V_0e^{-\gamma t}}{d} \\
    \Rightarrow B \cdot 2\pi r = \mu_0\sigma\pi r^2 E - \mu_0\epsilon \epsilon_0\gamma\pi r^2 E \\
    \Rightarrow \vec{B} = \frac{\mu_0 E r}{2} \left( \sigma - \epsilon \epsilon_0 \gamma \right) \hat{\phi}
\end{lefteqs}


לכל $b<r<a$ נסתכל על לולאה אמפירית ברדיוס $r$:
\begin{lefteqs}
    \oint_C \vec{B}\cdot d\vec{\ell} = 2\pi r B \\
    \mu_0 I_{in} = \mu_0 \cdot J \cdot S = \mu_0 \sigma E \pi b^2
\end{lefteqs}

זרם העתקה קיים גם בגליל:

\begin{lefteqs}
    I_{d,1}=-\gamma\epsilon\epsilon_0E\pi b^2,
\end{lefteqs}

וגם באוויר סביבו:
\begin{lefteqs}
    I_{d,2}=-\gamma\epsilon_0E\pi(r^2-b^2)
\end{lefteqs}

\begin{lefteqs}
    2\pi rB
    =\mu_0\sigma E\pi b^2-\mu_0 \gamma \epsilon \epsilon_0 E \pi b^2 - \mu_0 \gamma \epsilon_0 E \pi (r^2 - b^2) \\
    \vec{B} = \frac{\mu_0 E}{2r} \left( \sigma b^2 - \gamma \epsilon \epsilon_0 b^2 - \gamma \epsilon_0 (r^2 - b^2) \right) \hat{\phi}
\end{lefteqs}

\textbf{ג)} 

נמצא את וקטור פוינטינג בתוך המוליך ($0<r<b$):

ע"פ 4.92:

\begin{lefteqs}
    \vec{P}=\vec{E}\times\vec{H} =
    \begin{vmatrix}
        \hat{r} & \hat{\phi} & \hat{z}\\
        E_r & E_\phi & E_z\\
        H_r & H_\phi & H_z
    \end{vmatrix}
    = \begin{vmatrix}
        \hat{r} & \hat{\phi} & \hat{z}\\
        0 & 0 & E_z\\
        0 & H_\phi & 0
    \end{vmatrix} \\
    = (-E_z H_\phi)\hat{r} = (-E_z \cdot \frac{B_\phi}{\mu_0})\hat{r} \\
    = -\frac{V_0 \cdot e^{-\gamma t}}{d} \cdot \frac{\frac{1}{2} \cdot \frac{V_0 e^{-\gamma t}}{d}\cdot r(\sigma - \epsilon \epsilon_0 \gamma)}{\mu_0} \hat{r} \\
    \Rightarrow \vec{P} = -\frac{1}{2d} \cdot V_0^2 e^{-2\gamma t} \cdot r(\sigma -\epsilon \epsilon_0 \gamma) \hat{r}
\end{lefteqs}

\textbf{ד)} 
\begin{lefteqs}
   \vec{P} = -E \cdot \frac{E}{2r} (\sigma b^2 - \gamma \epsilon \epsilon_0 b^2 -\gamma \epsilon_0 (r^2 - b^2)) \hat{r} \\
   \vec{p} = \frac{-1}{2r} \cdot \frac{V_0^2 e^{-2\gamma t}}{d^2} \cdot (\sigma b^2 - \gamma \epsilon \epsilon_0 b^2 -\gamma \epsilon_0(r^2-b^2)) \hat{r}
\end{lefteqs}

\vspace{3em}
\hrule
\vspace{2em}

\textbf{\large שאלה 3}
נתון קבל עשוי משני לוחות מעגליים ברדיוס $a$ והמרחק ביניהם $b<<a$.

הקבל מחובר למקור מתח $V(t)=V_0\cos(\omega t)$.

\textbf{א)} 

ע"פ חוק אמפר:

\begin{lefteqs}
    \oint_C \vec{B} \cdot d\vec{\ell} = \mu_0 I_{in} + \mu_0 \epsilon_0 \frac{\partial \Phi_E}{\partial t} \\
    \vec{E} = \frac{V_0 \cos(\omega t)}{b} \hat{z} \Rightarrow E_z = \frac{V_0 \cos(\omega t)}{b} \\
    \underline{\forall r<a:} \\
    \Phi_E = E \cdot S = \frac{\pi r^2 V_0 \cos(\omega t)}{b} \\
    \partial \Phi_E / \partial t = \frac{-\pi r^2 \omega V_0 \sin(\omega t)}{b} \\
    I_{in} = 0 \\
    \vec{B} \cdot 2\pi r = \mu_0 \epsilon_0 \frac{-\pi r^2 \omega V_0 \sin(\omega t)}{b} \hat{\phi} \\
    \Rightarrow B_\phi(r,t) = -\frac{\mu_0 \epsilon_0 \omega r V_0 \sin(\omega t)}{2b}
\end{lefteqs}

\textbf{ב)} 

על מנת לחשב את $E_1$ כתוצאה מ $B_1$ נעזר בחוק פאראדיי האינטגרלי:

\begin{lefteqs}
    \oint_C \vec{E} \cdot d\vec{\ell} = -\frac{d}{dt} \iint_S \vec{B} \cdot \hat{n} dS \\
    B_\phi(r,t) = -\frac{\mu_0 \epsilon_0 \omega r V_0 \sin(\omega t)}{2b} \\
\end{lefteqs}

נעזר בלולאה אמפירית מלבנית, כאשר שני צדדיה נחים על שני לוחות הקבל, צד אחד נמצא במרכז שני הלוחות ($r=0$) והצד השני ברדיוס משתנה $r$:

\begin{lefteqs}
    \oint_C \vec{E} \cdot d\vec{\ell} = E(b+0+r-r) = bE \\
    -\frac{d}{dt} \iint_S \vec{B} \cdot \hat{n} dS = -\frac{d}{dt} \int_0^r B_\phi^1(r') b dr' = -\frac{d}{dt} \int_0^r -\frac{\mu_0 \epsilon_0 \omega r' V_0 \sin(\omega t)}{2b} b dr' \\
    = -\frac{d}{dt} \frac{b\mu_0 \epsilon_0 \omega V_0 \sin(\omega t)}{2b} \int_0^r r' dr' = -\frac{d}{dt} (\frac{\mu_0 r^2 \epsilon_0 \omega V_0 \sin(\omega t)}{4}) \\
    \Rightarrow bE = \frac{\mu_0 r^2 \epsilon_0 \omega^2 V_0 \cos(\omega t)}{4} \\
    \Rightarrow E_z^1 = \frac{-\mu_0 r^2 \epsilon_0 \omega^2 V_0 \cos(\omega t)}{4b}
\end{lefteqs}

\textbf{ג)} 
נחשב את $B_\phi^2 \, , E_z^2 \, , B_\phi^3 \, , E_z^3$

נעזר שוב בחוקי אמפר ופרדיי כדי למצוא את הערכים:

\begin{lefteqs}
    B_\phi \cdot 2\pi r = \mu_0 \epsilon_0 \frac{\partial \Phi_E^1}{\partial t} = \mu_0 \epsilon_0 \frac{\partial}{\partial  t} (E \cdot S) \\
    \Rightarrow B_\phi^n = \frac{\mu_0 \epsilon_0}{2\pi r} \cdot \frac{d}{dt} \int_0^r E_z^{n-1} 2\pi r' dr' \\
    B_\phi^2 = \frac{\mu_0 \epsilon_0}{r} \cdot \frac{d}{dt} \int_0^r \frac{\mu_0 r'^2 \epsilon_0 \omega^2 V_0 \cos(\omega t)}{4b} r' dr' \\
    B_\phi^2 = \frac{\mu_0 \epsilon_0}{r} \cdot \frac{\omega^2 \mu_0 \epsilon_0 V_0}{4b} \cdot \frac{d}{dt} \int_0^r -r'^3 \cos(\omega t) dr' \\
    = -\frac{(\mu_0 \epsilon_0)^2 \omega^2 V_0}{4b r} \cdot \frac{r^4}{4} \cdot (-\omega \sin(\omega t)) \\
    \Rightarrow B_\phi^2 =\frac{(\mu_0 \epsilon_0)^2 (\omega r)^3 V_0 \sin(\omega t)}{16b} \\
    E_z^2 = -\frac{1}{b} \frac{d}{dt} \int_0^r B_\phi^2(r') b dr' = \frac{(\mu_0 \epsilon_0)^2 (\omega r)^4 V_0 \cos(\omega t)}{64b} \\
    B_\phi^3 = \frac{\mu_0 \epsilon_0}{r} \cdot \frac{d}{dt} \int_0^r E_z^2(r') r' dr' = -\frac{(\mu_0 \epsilon_0)^3 (\omega r)^5 V_0 \sin(\omega t)}{384b} \\
    E_z^3 = -\frac{1}{b} \cdot \frac{d}{dt} \int_0^r B_\phi^3(r') b dr' = -\frac{(\mu_0 \epsilon_0)^3 (\omega r)^6 V_0 \cos(\omega t)}{2304b}
\end{lefteqs}

\textbf{ד)} 
\begin{lefteqs}
    E_z(r,t) = E_z^0 + E_z^1 + E_z^2 + E_z^3 \\
    E_z = \frac{V_0 \cos(\omega t)}{b} \left[1 - \frac{(\omega r / c)^2}{4} + \frac{(\omega r / c)^4}{64} - \frac{(\omega r / c)^6}{2304}\right] \\
    B_\phi = -\frac{V_0 \sin(\omega t)}{b c} \left[ \frac{\omega r / c}{2} - \frac{(\omega r / c)^3}{16} + \frac{(\omega r / c)^5}{384} \right]
\end{lefteqs}
כאשר $c = \frac{1}{\sqrt{\epsilon_0 \mu_0}}$ כפי שנכתב בהערה בעמוד 243

\textbf{ה)} 
נשווה לתוצאות מ 4.80 ו 4.81:
\begin{lefteqs}
    E(r,t) =  \frac{V_0}{d} \frac{J_0(kr)}{J_0(ka)} \cos(\omega t) \hat{z} \\
    B(r,t) = -\sqrt{\epsilon_0 \mu_0} \frac{V_0}{d} \frac{J_1(kr)}{J_0(ka)} \sin(\omega t) \hat{\phi}
\end{lefteqs}
כאשר פיתוחי מקלורן של פונקצית בסל הם:
\begin{lefteqs}
    J_0(x) = 1 - \frac{x^2}{4} + \frac{x^4}{64} - \frac{x^6}{2304} + \cdots \\
    J_1(x) = \frac{x}{2} - \frac{x^3}{16} + \frac{x^5}{384} - \cdots
\end{lefteqs}

בקירוב הקוואזיסטטי ($\omega \to 0 \, ,J_0(ka) \to 1$) גאשר $k=\frac{\omega}{c}$
\begin{lefteqs}
    E(r,t) = \frac{V_0 \cos(\omega t)}{d} \left[1 - \frac{(\omega r / c)^2}{4} + \frac{(\omega r / c)^4}{64} - \frac{(\omega r / c)^6}{2304}\right] \hat{z} \\
    B(r,t) = -\frac{V_0 \sin(\omega t)}{d c} \left[ \frac{\omega r / c}{2} - \frac{(\omega r / c)^3}{16} + \frac{(\omega r / c)^5}{384} \right] \hat{\phi}
\end{lefteqs}

קיבלנו תוצאה זהה עד כדי הקירוב שחושב


\vspace{3em}
\hrule
\vspace{2em}

\textbf{\large שאלה 4}

נתון הפונציאל הוקטורי מתאר גל אלקטרומגנטי $A=\frac{\alpha \sin(k(r-ct))}{r}  \hat{z}$

הקאורדינטה הכדורית $r$ מתארת את המרחק ממקור הגל

\textbf{א)} 
נחשב את השדה המגנטי בקאורדינטות גליליות

בקאורדינטות גליליות: $r=\sqrt{\rho^2 + z^2}$

\begin{lefteqs}
    \vec{B} = \vec{\nabla} \times \vec{A}
    \vec{B} = \frac{1}{\rho} \begin{vmatrix}
        \hat{\rho} & \hat{\phi} & \hat{z} \\
        \frac{\partial}{\partial \rho} & \frac{\partial}{\partial \phi} & \frac{\partial}{\partial z} \\
        A_\rho & \rho A_\phi & A_z
    \end{vmatrix}
    = \frac{1}{\rho} \cdot -\rho \hat{\phi} \cdot \frac{\partial A_z}{\partial \rho} \\
    = -\alpha \hat{\phi} \cdot \frac{k\rho \cdot \cos(k(\sqrt{\rho^2 + z^2} - ct)) - \frac{\rho \sin(k(\sqrt{\rho^2 + z^2} - ct))}{\sqrt{\rho^2 + z^2}}}{(\sqrt{\rho^2 + z^2})^2} \\
    \Rightarrow \vec{B} = -\alpha \rho \hat{\phi} \left( \frac{k\cos(k(r-ct))}{r^2} - \frac{\sin(k(r-ct))}{r^3} \right)
\end{lefteqs}

\textbf{ב)} 
בקאורדינטות כדוריות $\rho=r\sin(\theta)$
\begin{lefteqs}
    \vec{B} = -\alpha r \sin(\theta) \hat{\phi} \left( \frac{k\cos(k(r-ct))}{r^2} - \frac{\sin(k(r-ct))}{r^3} \right) \\
    \Vec{B} = -\alpha \sin(\theta) \hat{\phi} \left( \frac{k\cos(k(r-ct))}{r} - \frac{\sin(k(r-ct))}{r^2} \right)
\end{lefteqs}

\textbf{ג)} 
\begin{lefteqs}
    \vec{\nabla} \times \vec{B} = \frac{1}{r^2 \sin(\theta)} \begin{vmatrix}
        \hat{r} & r\hat{\theta} & r\sin(\theta)\hat{\phi} \\
        \frac{\partial}{\partial r} & \frac{\partial}{\partial \theta} & \frac{\partial}{\partial \phi} \\
        B_r & rB_\theta & r\sin(\theta)B_\phi
    \end{vmatrix} \\
    = \frac{1}{r^2\sin(\theta)}(\hat{r} \cdot \frac{\partial}{\partial \theta}(r\sin(\theta)B_\phi) - r\hat{\theta} \cdot \frac{\partial}{\partial r}(r\sin(\theta)B_\phi)) \\
    = \frac{1}{r^2\sin(\theta)} \left( r\hat{r} \cdot 2\cos(\theta) \cdot B_\phi - \hat{\theta} \cdot r\sin{\theta} \cdot (B_\phi + r\cdot\frac{\partial B_\phi}{\partial r}) \right) \\
    \frac{\partial B_\phi}{\partial r} = -\alpha \sin(\theta) \left(k \frac{-k\sin(k(r-ct))\cdot r - \cos(k(r-ct))}{r^2} - \frac{k\cos(k(r-ct))r^2 - 2r\sin(k(r-ct))}{r^4} \right) \\
    \beta = k(r-ct) \\
    \frac{\partial B_\phi}{\partial r} = -\alpha \sin(\theta) ( -\frac{k^2\sin(\beta)}{r} - \frac{k\cos(\beta)}{r^2} - \frac{k\cos(\beta)}{r^2} + \frac{2\sin(\beta)}{r^3} ) \\
    = \alpha \sin(\theta) \left( \frac{k^2\sin(\beta)}{r} + \frac{2k\cos(\beta)}{r^2} - \frac{2\sin(\beta)}{r^3} \right) \\
    B_\phi + r\frac{\partial B_\phi}{\partial r} = \alpha \sin(\theta) (k^2 \sin(\beta) + \frac{2k\cos(\beta)}{r} - \frac{2\sin(\beta)}{r^2} - \frac{k\cos(\beta)}{r} + \frac{sin(\beta)}{r^2}) \\
    = \alpha \sin(\theta) (k^2 \sin(\beta) + \frac{k\cos(\beta)}{r} - \frac{\sin(\beta)}{r^2}) \\
    \vec{\nabla} \times \vec{B} = \hat{r} (\frac{1}{\sin(\theta)} \cdot 2\cos(\theta) (-\alpha) \sin(\theta)(\frac{k\cos(\beta)}{r} - \frac{\sin(\beta)}{r^2})  \\
    \qquad \qquad -\hat{\theta} (\frac{1}{r} \alpha \sin(\theta) (k^2 \sin(\beta) + \frac{k\cos(\beta)}{r} - \frac{\sin(\beta)}{r^2}))))
\end{lefteqs}

מכיוון ש $J=0$ ע"פ חוק אמפר המתוקן: $\vec{\nabla} \times \vec{B} = \mu_0 \epsilon_0 \frac{\partial E}{\partial t}$

\begin{lefteqs}
    E = c^2 \int \vec{\nabla} \times \vec{B} dt \\
    E_r = c^2 \cdot \frac{-2 \alpha \cos(\theta)}{r} \int \frac{k\cos(k(r-ct))}{r} - \frac{\sin(k(r-ct))}{r^2} dt \\
    = -\frac{-2\alpha c^2 \cos(\theta)}{r} \left( \frac{-\sin(\beta)}{rc} - \frac{\cos(\beta)}{r^2 kc} \right) \\
    E_\theta = c^2 \cdot \frac{-\alpha \sin(\theta)}{r} \int (k^2 \sin(\beta) + \frac{\cos(\beta)}{r} - \frac{\sin(\beta)}{r^2}) dt \\
    = \frac{\alpha c \sin(\theta)}{r} (-k\cos(\beta) + \frac{\sin(\beta)}{r} + \frac{\cos(\beta)}{k r^2}) \\
    \vec{E} = E_r \hat{r} + E_\theta \hat{\theta} \, , \beta = k(r-ct)
\end{lefteqs}


\vspace{3em}
\hrule
\vspace{2em}

\textbf{\large שאלה 5}

\textbf{א.1)} 
נבדוק את התקיימות חוק פרדיי $\vec{\nabla} \times \vec{E} = -\frac{\partial \vec{B}}{\partial t}$

\begin{lefteqs}
    -\frac{\partial \vec{B}}{\partial t} = -\frac{\partial}{\partial t} \left( \alpha \sin(\theta)(\frac{k\cos(k(r-ct))}{r} - \frac{\sin(k(r-ct))}{r^2}) \right) \\
    = \alpha \sin(\theta) \left( \frac{k \cdot (-kc) \cdot -\sin(\beta)}{r} - \frac{-kc \cdot \cos(\beta)}{r^2} \right) \\
    = kc \alpha \sin(\theta) (\frac{k\sin(\beta)}{r} + \frac{\cos(\beta)}{r^2}) \\
    \vec{\nabla} \times \vec{E} = \frac{1}{r^2 \sin(\theta)} \begin{vmatrix}
        \hat{r} & r\hat{\theta} & r\sin(\theta)\hat{\phi} \\
        \frac{\partial}{\partial r} & \frac{\partial}{\partial \theta} & \frac{\partial}{\partial \phi} \\
        E_r & rE_\theta & 0
    \end{vmatrix} \\
    = \frac{1}{r^2 \sin(\theta)} \left( r\sin(\theta) \hat{\phi} (\frac{\partial(rE_\theta)}{\partial r} - \frac{\partial E_r}{\partial \theta}) \right) \\
    = \hat{\phi} \cdot \frac{1}{r} [\frac{\partial}{\partial r}(\alpha c \sin(\theta)(-k \cos(\beta) + \frac{\sin(\beta)}{r} + \frac{\cos(\beta)}{kr^2})  \\
    - \frac{\partial}{\partial \theta}(\frac{2\alpha \cos(\theta) \cdot c}{r^2}(\sin(\beta) - \frac{\cos(\beta)}{kr})) ] \\
    = \hat{\phi} \frac{\alpha c}{r} [\sin(\theta)(k^2 \sin(\beta) + \frac{k\cos(\beta)}{r} - \frac{\sin(\beta)}{r^2} - \frac{\sin(\beta)}{r^2} - \frac{2\cos(\beta)}{kr^3}) + \sin(\theta)(\frac{2\sin(\beta)}{r^2} + \frac{2\cos(\beta)}{kr^3}))] \\
    = \hat{\phi} \frac{\alpha c \sin(\theta)}{r} (k^2 \sin(\beta) + \frac{k\cos(\beta)}{r} - \frac{2\sin(\beta)}{r} - \frac{2\cos(\beta)}{kr^2}\\
    \qquad + 2\frac{\sin(\beta)}{r^2} + \frac{2\cos(\beta)}{kr^3}) \\
    = \hat{\phi} = k\alpha c \sin(\theta) (\frac{k\sin(\beta)}{r} + \frac{\cos(\beta)}{r^2}) \\
    \Rightarrow \vec{\nabla} \times \vec{E} = -\frac{\partial \vec{B}}{\partial t}
\end{lefteqs}

\textbf{א.2)} 
חוק גאוס בריק הוא $\vec{\nabla} \cdot \vec{E} = 0$
\begin{lefteqs}
    \vec{\nabla} \cdot \vec{E} = \frac{\partial E_r}{\partial r} + \frac{2}{r} E_r + \frac{1}{r} (\frac{\partial E_\theta}{\partial \theta} + \frac{\cos(\theta)}{\sin(\theta)} E_\theta) \\
    2 \alpha \cos(\theta) c \left[ \frac{k\cos(\beta)}{r^2} -\frac{2\sin(\beta)}{r^3} - \frac{\sin(\beta)}{r^3} - \frac{3\cos(\beta)}{kr^4} \right] \\
    + \frac{4\alpha \cos(\theta) c}{r^3} (\sin(\beta) + \frac{\cos(\beta)}{kr}) + \frac{2c \cos(\theta)}{r^2}(-k\cos(\beta) + \frac{\sin(\beta)}{r} + \frac{\cos(\beta)}{kr^2}) \\
    + \frac{2c\cos(\theta)}{r^2}(-k\cos(\beta) + \frac{\sin(\beta)}{r} + \frac{\cos(\beta)}{kr^4}) \\
    = \alpha c \cos(\theta) [\frac{-6\sin(\beta)}{r^3} - \frac{6\cos(\beta)}{kr^4} + \frac{2k\cos(\beta)}{r^2} + \frac{4\sin(\beta)}{r^3} + \frac{4\cos(\beta)}{kr^4} - \frac{k\cos(\beta)}{r^2} + \frac{\sin(\beta)}{r^3}\\
    \qquad + \frac{\cos(\beta)}{kr^4} - \frac{k\cos(\beta)}{r^2} + \frac{\sin(\beta)}{r^3} + \frac{\cos(\beta)}{kr^4}) ] \\
    \frac{\sin(\beta)}{r^3}(4+1+1-6) + \frac{k\cos(\beta)}{r^2}(2-1-1) + \frac{\cos(\beta)}{kr^4}(1+4+1-6) = 0 \\
    \Rightarrow \vec{\nabla} \cdot \vec{E} = 0
\end{lefteqs}

\textbf{ב)} 
לא.

ידועה הזהות של $\vec{\nabla} \cdot (\vec{\nabla} \times \vec{A})=0$

מכיוון שידוע $\vec{B} = \vec{\nabla} \times \vec{A}$ אז

\begin{lefteqs}
    \vec{\nabla} \cdot \vec{B} = \vec{\nabla} \cdot (\vec{\nabla} \times \vec{A}) = 0
\end{lefteqs}

\end{document}
